% ============================================================================
%  04-06: ANÁLISIS CRÍTICO Y OBJETIVOS DE DESARROLLO SOSTENIBLE (ODS)
% ============================================================================

\section{An\'alisis cr\'itico y alineaci\'on con los ODS}

\subsection{An\'alisis cr\'itico de las decisiones de dise\~no}

\subsubsection{Arquitectura multiesquema frente a esquema \'unico}

Dividir el modelo en cuatro esquemas de negocio (\texttt{catalogo}, \texttt{ventas}, \texttt{operaciones} y \texttt{auditoria}) manteniendo ASP.NET Core Identity en \texttt{public} introduce una configuraci\'on expl\'icita en Fluent API, pero resuelve problemas estructurales reales:

\begin{itemize}
  \item \textbf{Compatibilidad sin interrupci\'on de autenticaci\'on:} Conservar \texttt{AspNetUsers} y \texttt{AspNetRoles} en \texttt{public}, eliminando la columna redundante \texttt{FullName} e incorporando las claves for\'aneas hacia \texttt{operaciones.Sucursales} y \texttt{operaciones.Regiones}, permiti\'o migrar la base de datos en Supabase sin invalidar las cuentas ni las sesiones existentes.
  \item \textbf{Principio de menor privilegio:} Permite otorgar permisos \texttt{USAGE} y \texttt{SELECT/INSERT} diferenciados por esquema seg\'un el rol operativo, aislando las tablas inmutables de \texttt{auditoria} del acceso directo de los m\'odulos comerciales \citep{postgresql2024}.
  \item \textbf{Mantenibilidad modular:} Cada dominio funcional agrupa sus propias tablas, vistas y rutinas (\texttt{ventas.sp\_confirmar\_venta}, \texttt{catalogo.fn\_recalcular\_stock\_combo}), facilitando futuras ampliaciones.
\end{itemize}

\subsubsection{Arco exclusivo en \texttt{PedidoItems} frente a referencia polim\'orfica}

En el estado inicial, \texttt{PedidoItems} usaba el par \texttt{(Tipo, ProductoId)} sin clave for\'anea.
Sustituirlo por las dos claves for\'aneas anulables \texttt{ComponenteId} y \texttt{ComboId} gobernadas por la restricci\'on \texttt{ck\_pi\_exclusivo} introduce un campo nulo por fila, pero garantiza a nivel del motor PostgreSQL que ning\'un componente o combo con historial de ventas pueda ser eliminado f\'isicamente (\texttt{ON DELETE RESTRICT}) ni referenciado si no existe \citep{date2004, silberschatz2019}.

\subsubsection{Uso deliberado de \texttt{JSONB} en \texttt{CarritosAbandonados}}

Mantener \texttt{ItemsJson} como \texttt{jsonb} indexado con \texttt{GIN} en lugar de crear una tabla detalle adicional responde a la naturaleza vol\'atil del carrito previo al pago: permite persistir instant\'aneas r\'apidas desde la sesi\'on web, consultar su contenido con el operador \texttt{@>} en milisegundos y delegar la validaci\'on estricta en 3FN al momento del cierre transaccional en \texttt{ventas.Pedidos} y \texttt{ventas.PedidoItems}.

\subsection{Alineaci\'on con los Objetivos de Desarrollo Sostenible}

\subsubsection{ODS 8 --- Trabajo decente y crecimiento econ\'omico}

El redise\~no implementado en TechParts ERP respalda las metas 8.2 y 8.3 de modernizaci\'on tecnol\'ogica y formalizaci\'on de las PYME \citep{ods8_2015}:

\begin{enumerate}
  \item \textbf{Productividad operativa y decisiones basadas en datos:} Las vistas anal\'iticas (\texttt{vw\_ventas\_mensuales}, \texttt{vw\_combos\_rentabilidad}, \texttt{vw\_clientes\_activos}) y los siete \'indices desplegados permiten obtener indicadores de facturaci\'on, m\'argenes y tiempos de entrega por sucursal sin procesos manuales en hojas de c\'alculo.
  \item \textbf{Escalabilidad multiregional:} El esquema \texttt{operaciones} (\texttt{Regiones}, \texttt{Sucursales}, \texttt{Proveedores}) soporta la expansi\'on comercial entre Cochabamba, La Paz y Santa Cruz con trazabilidad de comisiones de venta.
\end{enumerate}

\subsubsection{ODS 12 --- Producci\'on y consumo responsables}

En el sector de hardware electr\'onico, la gesti\'on deficiente de inventarios incrementa la generaci\'on de residuos de aparatos el\'ectricos y electr\'onicos (RAEE).
La arquitectura desplegada materializa el ODS~12 directamente en el modelo de datos \citep{ods12_2015}:

\begin{enumerate}
  \item \textbf{Prevenci\'on de sobreventa y compras innecesarias:} Las restricciones \texttt{CHECK (Stock >= 0)}, el umbral \texttt{StockMinimo} por producto y el bloqueo pesimista en \texttt{sp\_confirmar\_venta} evitan quiebres artificiales de inventario y env\'ios fallidos por inconsistencias de stock.
  \item \textbf{Ciclo de vida y trazabilidad RAEE (\texttt{MermasInventario}):} El atributo \texttt{EstadoObsolescencia} en \texttt{catalogo.Componentes} (\texttt{Vigente}, \texttt{PorDescontinuar}, \texttt{Obsoleto}, \texttt{Merma}) junto con la tabla \texttt{auditoria.MermasInventario} exigen registrar la disposici\'on final verificable de cada baja (\texttt{ReciclajeCertificado}, \texttt{Reacondicionado}, \texttt{Donacion} o \texttt{BajaDefinitiva}), como se demostr\'o en la Consulta~9 con el seguimiento de unidades reacondicionadas y enviadas a reciclaje certificado RAEE.
\end{enumerate}
